\chapter{任务队列与分布式边界}\label{ch:scale}
\begin{tcolorbox}[title=章首知识地图,colframe=Ink]
待访问队列 $\longrightarrow$ 工作者 $\longrightarrow$ 去重与写入 $\longrightarrow$ 失败重投与监控。
\end{tcolorbox}

\begin{tcolorbox}[title=学习目标与前置知识,colframe=Ink!65,colback=Ink!2]
前置：去重键与幂等写入。目标：说明重复投递如何在存储层被消解。\par
输入：含重复项的任务队列；输出：三条唯一记录。
\end{tcolorbox}

\section{为什么需要队列}
当页面数量增加，待访问地址需要持久保存、分配给工作者并记录处理状态。队列把“发现新地址”和“处理地址”分开；多个工作者可以分担任务，但同时带来重复投递、顺序变化和故障恢复问题。

\begin{definition}{至少一次处理与幂等}
消息系统为了避免丢任务，可能把同一地址交给工作者不止一次。若重复处理不会产生重复结果，就称写入具有幂等性。分布式系统中不能只靠每个进程自己的 \texttt{set} 去重；唯一键或协调存储应提供共同约束。
\end{definition}

\section{单进程模拟重复投递}
示例用线程和队列模拟多个工作者，队列里故意放入重复任务。它只演示“重复投递必须预期”这一点，并不是分布式爬虫；真正的跨机器队列还要处理确认、超时与持久化。

\begin{example}{重复任务只计一次}{基础}
\VerbatimInput{examples/ch21_queue.py}
\end{example}

\begin{method}{从单机扩展的顺序}
先固定字段契约和单页解析测试，再让单机分页稳定运行；接着引入持久化队列与唯一键，最后才增加工作者。每扩一层，都记录吞吐量、失败率和站点负载。机器多了，错误也会组团出行。
\end{method}

\section{选读：智能抽取的适用位置}
规则选择器适合结构稳定、字段明确的页面；结构多变时可以考虑模板分类、启发式或模型辅助提取。但“看起来像标题”仍需来源、字段类型与样本复核。智能抽取应负责提出候选，验收仍由明确的数据契约完成。

\begin{tcolorbox}[title=本章验收,colframe=Teal,colback=Teal!3]
运行单进程队列示例：同一次运行中重复任务只计一次，得到编号 1、2、3。跨次运行的持久去重由第\ref{ch:project}章的文件数据库验证。
\end{tcolorbox}

\section{练习}
\begin{enumerate}[label=\arabic*.]
\item 把同一任务放入队列三次，说明为何结果仍应只计一次。
\item 解释进程内去重集合为何不能直接扩展到两台机器。
\item 为模型辅助提取设计一项人工抽样复核指标。
\end{enumerate}

\section{补充：BloomFilter 与可恢复任务队列}\label{rev:crawler:redis}
BloomFilter 用位数组和多个哈希函数判断成员关系。在没有删除及实现错误等前提下，回答“不存在”可作为不存在证据，回答“存在”可能是假阳性。插入 $n$ 个元素、位数 $m$、哈希数 $k$ 时，常用近似误报率为 $(1-e^{-kn/m})^k$，较优 $k\approx(m/n)\log2$。不能仅凭“可能存在”就永久丢弃必须完整采集的任务。
\begin{Verbatim}
# Redis 需支持 BF.*；以下是 redis-cli 终端命令。
BF.RESERVE notes:demo:bloom 0.001 100000
BF.ADD notes:demo:bloom fixture:1
BF.EXISTS notes:demo:bloom fixture:1
\end{Verbatim}
可用 BloomFilter 降低精确集合的查询量，最终完整性仍由精确键与幂等结果存储保证。BF.RESERVE 的容量是设计参数，不是数据库行数承诺；服务端是否提供 BF.* 应先检查，不能认为每个旧版 Redis 都内置。

\subsection{结果提交后确认，崩溃后重新接管}
Redis Streams 消费组维护待确认列表。工作者读取后，先完成幂等结果写入并提交，再 XACK；若在提交后、确认前崩溃，重投会发生，但唯一键防止重复结果。XAUTOCLAIM 可接管空闲时间超过阈值的待确认消息；阈值要高于正常处理时间，长任务还需续租或其他协调，防止两名工作者同时处理。
\VerbatimInput{examples/revision/redis_resume.py}
这是可连接本机 Redis 6.2+ 的教学骨架：先执行 seed，再执行 work，三条消息最终对应两条唯一结果。示例每次只取一批，新任务较多时要循环 work。SQLite 文件只演示同机持久化结果；跨机器应改用共享的事务数据库及唯一键，不能让每台机器各存一份 SQLite 后声称全局去重。
发生异常的消息保持未确认，达到空闲阈值后可重取；正式部署还需重试上限、死信队列、监控、Redis 持久化与恢复演练，且不能随意裁剪尚未处理的 Stream。该骨架不提供恰好一次的跨系统事务保证。本次做了源码检查，Redis 服务联调步骤在随书验收说明中单列。
概念与接管语义参见 \href{https://redis.io/docs/latest/develop/data-types/probabilistic/bloom-filter/}{Redis BloomFilter} 与 \href{https://redis.io/docs/latest/commands/xautoclaim/}{XAUTOCLAIM 官方文档}。

